% ---------------------------------------------------------------------------
% riccati-cert-lean/main.tex -- the VERIFICATION COMPANION NOTE.
%
% CONTENT CONTRACT: verify/plan-map.json (23 planned claims, PLAN-L-001..023)
% for the original Lean-corpus sections, plus the 2026-08-18 widening (Adi
% directive, papers/PUBLISHING-PROMPTS.md Trigger): the companion also carries
% the executable-receipts catalogue and the PodiumNote.lean mechanization of
% the PUBLISHED technical note's theorems. Claims of the two new sections mint
% at extraction (LCLAIM-024+) and are anchored by summarizes/context pointers
% into the full paper's register or by summarizes-allowlist.json entries.
% Section labels:
%   sec:what           N1  (PLAN-L-001, -002; opening widened 2026-08-18)
%   sec:corpus         N2  (PLAN-L-003, -004, -005, -006)
%   sec:acceptance     N3  (PLAN-L-007, -008, -009)
%   sec:bridge         N3.1 subsection of N3 (PLAN-L-010, -011)
%   sec:soundness      N4  (PLAN-L-012, -013, -014, -015, -016)
%   sec:wdd            N5  (PLAN-L-017)
%   sec:receipts       N6  (2026-08-18 widening: the executable receipts)
%   sec:published-note N7  (2026-08-18 widening: PodiumNote.lean)
%   sec:scope          N8  (PLAN-L-018..-023)
% Theorem statements are given in prose + math with exact Lean identifiers in
% \texttt{}; no transliterated Lean listings (the source uses Unicode syntax
% that the pinned pdflatex/Computer Modern stack does not carry, and an
% ASCII-transliterated quotation would not be the corpus's actual text).
% Lean-core tokens that are not corpus declarations (#print axioms, propext,
% Classical.choice, Quot.sound, sorryAx) are set with \verb so check_lean.py's
% cited-identifier scan ranges over corpus declarations only.
% Preamble conventions carried from papers/riccati-cert-simple/main.tex so the
% note matches the paper family; title and author line per the simplified
% paper's leannote bibliography entry and author block.
% ---------------------------------------------------------------------------
\documentclass[11pt]{article}
\usepackage[margin=1in]{geometry}
\usepackage{amsmath,amssymb,amsthm}
\usepackage{booktabs}
\usepackage[hidelinks]{hyperref}

\newtheorem{theorem}{Theorem}
\newtheorem{lemma}{Lemma}
\theoremstyle{definition}
\newtheorem{assumption}{Assumption}
\newtheorem{definition}{Definition}
\newtheorem{example}{Example}
\theoremstyle{remark}
\newtheorem{remark}{Remark}
\newcommand{\Q}{\mathbb{Q}}
\newcommand{\R}{\mathbb{R}}
\DeclareMathOperator{\lcm}{lcm}
\DeclareMathOperator{\den}{den}

% Typography: long typewriter identifiers (border\_band\_psd,
% podium.verify.riccati) carry no break points, so a line ending in one
% overflows the measure. Allow a break after each underscore, and give TeX a
% third-pass stretch budget so an unbreakable token moves to the next line
% instead of running past it. check_refs fails the build on any overfull box
% wider than 2pt, so this must hold.
% OT1's \_ is a drawn RULE, not a character, so every identifier extracted
% from the PDF lost its underscores (border_band_psd came out "border band
% psd") -- copy-paste and any text indexer saw the wrong name. Emit the
% typewriter font's real underscore glyph instead and let pdftex map it to
% U+005F. The glyph exists only in the typewriter font -- char 95 of cmr is not
% an underscore -- so a roman \_ keeps the rule. The fontdimen3 test is the
% standard "is this monospace" check: tt fonts have no interword stretch.
\input{glyphtounicode}
\pdfgentounicode=1
\renewcommand{\_}{\ifdim\fontdimen3\font=0pt\relax\char`\_\else\textunderscore\fi\hspace{0pt}}
\emergencystretch=3em

% Draft version. SINGLE SOURCE OF TRUTH for the version number: papers/build.sh
% names the output <paper-dir>-v<paperversion>.pdf from this macro, and
% verify/_common.py parses it to find the build products. Bump it here and
% nowhere else; the title page carries it so a loose PDF names its version.
\newcommand{\paperversion}{1}

\title{Machine-Checked Foundations for the Exact-Rational Riccati Certificate:\\
A Lean Companion Note}
\author{Adi Oltean\thanks{Starcloud, Inc. \texttt{adi@starcloud.com}. Working draft.
AI coding assistants were used for drafting, for the reference implementation and
verification scripts, and for adversarial review; all results were checked by the author.}}
\date{Version \paperversion{} --- August 2026}

\begin{document}
\maketitle

\section{What this note is}\label{sec:what}
This note is the verification companion to the paper on exact-rational
certification of block-banded trajectory QCQPs~\cite{riccatipaper}, and it
does three jobs. It documents the Lean~4 formalization behind the
machine-checked results cited there: each such citation resolves to the
numbered section here that states the corresponding declaration and its
hypotheses. It catalogues the executable receipts behind the paper's
measured and instance-level results (\S\ref{sec:receipts}). It records the
mechanization of the published Podium technical note's theorems
(\S\ref{sec:published-note}).

The corpus machine-checks classical facts along the paper's trusted path
rather than new mathematics: what it supplies is a tolerance-free,
mechanically verified chain, not novelty in the individual lemmas.

\section{The corpus and its pin}\label{sec:corpus}
The trusted path is mechanized in eighteen top-level Lean~4 files, all
\texttt{sorry}-free, in the Lean corpus of the Podium project
(\texttt{github.com/adi-oltean/podium}). Every declaration this note cites is
named explicitly, so a reader can locate it by name in any copy of the corpus.
The formalization is also available from the author at
\texttt{adi@starcloud.com}.

Every file type-checks against a pinned project built from
\texttt{leanprover/lean4:v4.33.0} and Mathlib at revision \texttt{v4.33.0}.
A per-file regression script drives the build: it reports pass or fail for
each file and surfaces \verb|#print axioms| for the headline declarations.
The axiom audit admits only \verb|propext|, \verb|Classical.choice| and
\verb|Quot.sound|, and never \verb|sorryAx|.

\section{The acceptance predicate, both directions}\label{sec:acceptance}
\texttt{ldlt\_isPSD} (\texttt{LDLTSound.lean}): if
$M=L^\top\,\mathrm{diagonal}(d)\,L$ and every pivot $d_i\ge0$, then $M$ is PSD
in the quantifier sense $\forall v,\ 0\le v^\top Mv$. This is the acceptance
direction; it holds at any order and over any linear ordered field, by the
expansion $v^\top Mv=\sum_i d_i(Lv)_i^2$.

\texttt{ldlt\_complete} (\texttt{SchurStep.lean}): a symmetric PSD matrix of
any order $m$ admits weights $d_j\ge0$ and vectors $v_j$ with
$M=\sum_j d_jv_jv_j^\top$. The induction peels one nonnegative-weighted
rank-one term per pivot, and it passes through singular pivots because a zero
pivot forces a zero border (\texttt{headBordered\_border\_zero}). This is the
converse of \texttt{ldlt\_isPSD} up to padding the factor to square; the
residual difference is purely dimensional ($k\le m$).

\texttt{isPSD\_iff\_ldlt} (same file) packages the two directions as one
equivalence: for symmetric $M$, $M$ is PSD if and only if such a
nonnegative-weight decomposition exists. This is the machine-checked form of
accept $\Leftrightarrow$ PSD, with \emph{accept} being exactly that
existential.

\subsection{The bridge to the library-standard notion}\label{sec:bridge}
\texttt{posSemidef\_iff\_isPSD} (\texttt{PosSemidefBridge.lean}): for a
rational $n\times n$ matrix $M$, Mathlib's \texttt{Matrix.PosSemidef}\,$M$
holds exactly when $M^\top=M$ and $\texttt{IsPSD}\,M$, the quantifier form
$\forall v,\ 0\le v^\top Mv$. Symmetry is derived by the equivalence, not
assumed of $M$. So the quantifier form is the library-standard notion
rather than a bespoke one. That bridge is compiled on its own and is not
linked inside Lean to the predicate the acceptance lemmas above are stated
against; the condition is stated in \S\ref{sec:scope}.

\section{Certificate soundness and closure}\label{sec:soundness}
\texttt{general\_sprocedure\_sound} (\texttt{GeneralLifting.lean}) takes
three hypotheses:
\begin{enumerate}
\item the bordered matrix built from Hessian $P$, half-linear part $b$ and
corner $c$ is PSD;
\item every multiplier $\lambda_i\ge0$;
\item the form identity: for every point, the bordered quadratic form at the
lift equals $f_0-\sum_i\lambda_if_i-t$.
\end{enumerate}
From these it concludes $t\le f_0(y)$ at every $y$ with all $f_i(y)\ge0$.
The theorem holds at any order $n$, at any number $m$ of constraints, and
over any linear ordered field.

Of those three, the form identity is the one not mechanized in general. The
abstract lifting identity \texttt{bordered\_lifting} (same file) is proved at
general $n$: the bordered matrix's quadratic form at $(x,1)$ equals
$x^\top Px+2b^\top x+c$. Its instantiation to a specific S-procedure assembly
is discharged per instance by \texttt{ring}, and the capstone theorems carry
that instantiation as a hypothesis too.

\texttt{certified\_global\_optimum} (\texttt{ClosureOptimum.lean}): given
multipliers $\lambda\ge0$, a witness $\bar x$ that minimizes the Lagrangian
$f_0-\sum_i\lambda_if_i$ over all points, and complementary slackness
$\sum_i\lambda_if_i(\bar x)=0$, the witness satisfies $f_0(\bar x)\le f_0(y)$
at every feasible $y$ --- the attainment leg, over any linear ordered field.
Its companion \texttt{bracket\_exact} (same file) adds feasibility of $\bar x$
and a supplied lower bound --- that $t\le f_0$ on the feasible set --- and
concludes both $t\le f_0(\bar x)$ and global optimality of $\bar x$. The
equality $t^\star=f_0(x^\star)$ that closes the bracket is by construction of
$t^\star$ and is not separately mechanized.

\texttt{certified\_optimum\_pipeline} (\texttt{Certificate.lean}) assembles
the chain as one theorem. Its hypotheses are an accepted exact factorization
$M=L^\top DL$ with pivots $\ge0$, multipliers $\lambda\ge0$, the form
identity, and a feasible Lagrangian-minimizing witness $\bar x$ with
complementary slackness. It returns three conclusions: $t$ lower-bounds
$f_0$ on the feasible set; $\bar x$ is a global optimum; and
$t\le f_0(\bar x)$. From these, $f_0(\bar x)=J^\star$ follows because
$\bar x$ is feasible and minimal. Two fully worked exact-$\Q$ instances, a
one-dimensional keep-out and a two-dimensional disk, run through the
verdict-level twin \texttt{certified\_optimum\_of\_verdict}, which starts
from the PSD verdict rather than from a factorization.

\section{The diagonal-dominance leg}\label{sec:wdd}
\texttt{isPSD\_of\_wdd} (\texttt{GershgorinPSD.lean}): symmetry together with
row dominance $\sum_{j\ne i}|M_{ij}|\le M_{ii}$ implies PSD at general $n$
over any linear ordered field, and the declaration takes exactly those two
hypotheses. This is the Gershgorin leg of the paper's diagonal-dominance
lemma.

\section{The executable receipts}\label{sec:receipts}
The paper's measured and instance-level results rest on a set of executable
receipts: exact-arithmetic Python scripts that construct the instances, run
them through the trusted checkers, and assert what the paper reports. This
section catalogues those receipts by role. Every receipt below ships in the
paper's supplementary bundle, together with the supporting modules the
drivers import, with one exception: \texttt{test\_riccati} is part of the
Podium library's own test suite. The bundle's own \texttt{README.md} gives the
exact setup and invocation for every receipt below: the library commit to
clone, the dependencies, and the command line for each script and for the
pytest suite. Throughout the receipts, the
bit-size of a rational is
\texttt{numerator.bit\_length()}${}+{}$\texttt{denominator.bit\_length()}.

The trusted checkers are not bundle contents: they live in the Podium
repository at \texttt{github.com/adi-oltean/podium}, and the receipts import
them from there (the library import requires NumPy). The band checkers are
\texttt{border\_band\_psd} and \texttt{block\_tridiag\_psd} in
\texttt{podium.verify.riccati}; retrieving them together with
\texttt{test\_riccati} requires commit \texttt{c96fcb5} of that repository. The archived Zenodo
release of the library, v0.8.1 (DOI 10.5281/zenodo.21302775), predates
\texttt{podium.verify.riccati}: it carries \texttt{podium.verify.bracket}
(the dense $M$-test) and \texttt{podium.verify.barrier}, but not the band
sweeps.

\paragraph{The family drivers.} Three drivers construct the families of the
paper's results table in exact \texttt{Fraction} arithmetic, and each
verifies its emitted certificate through the exact dense $M$-test. The
trusted checkers refuse floats; the CW state-transition matrix alone is
synthesized in floating point and rationalized before any certification.
\begin{description}
\item[\texttt{shared\_state\_exact.py}] constructs the scalar shared-state
family; it pins the results table's scalar row.
\item[\texttt{vector\_state.py}] constructs the vector shared-state family
($d=1..4$) and pins the table's vector row. Its probe \texttt{gap\_probe}
($\lambda_k>w_k$ at $N{=}4$, $d{=}2$) is the aggregate-PSD-violated instance
of the table's refusal row, on which the checker returns no certificate.
\item[\texttt{rpod\_cw.py}] constructs the CW-coupled keep-out family of the
paper's demonstration section and pins the table's CW row. The family's
data are mean motion $n_0=10^{-3}$\,rad/s, sampling phase
$\tau=n_0\,\Delta t=1/2$, weights $R=I_4$ and $W_k=w_kI_4$, per-stage radii
cycling $\{1/2,1,3/2\}$, and STM denominators capped at $10^5$. The driver
also prints the block pivots $S_k$, which grow measurably linearly in the
horizon: $780$ bits at $N{=}2$ to $17889$ at $N{=}24$.
\end{description}

\paragraph{The measurement engines.}
\begin{description}
\item[\texttt{cert\_metrics.py}] recomputes both bit columns of the results
table for the three closing families at the headline sizes, and re-verifies
each of those rows through all three trusted routes: the dense $M$-test, the
band-$H$ sweep, and the bordered-$M$ sweep.
\item[\texttt{cw\_anatomy.py}] recomputes, at all $23$ horizons
$N=2,\dots,24$, every denominator statistic the paper quotes for the
CW-coupled family, together with the closure verdicts and both bit columns.
It asserts the divisibility law behind the certified value's denominator,
and the value's size inequalities, at every horizon. It measures the
entrywise rationalization error of the CW state-transition matrix. Its
bordered-sweep section pins the corner-carrying sweep pivot by pivot against
the dense $LDL^\top$ of $M$ at the shortest horizons, which is the
test-pinning leg of the paper's bordered-extension remark.
\end{description}

\paragraph{The exact instance receipts.} Four receipts certify the paper's
named instances entirely over $\Q$; the fourth certifies two sibling
instances in one run.
\begin{description}
\item[\texttt{single\_keepout\_coupling\_gap.py}] is the paper's certified
duality-gap example. It asserts $J^\star=212/13$ by exact enumeration of the
branch-KKT candidates, certifies through a rational moment point that every
level-one certificate is capped at $815/64$, and certifies the level-one gap
$\ge2973/832$. The same receipt discharges the diagnosis: stationarity
recovers $\lambda^\star=24/13$, and $H(\lambda^\star)$ has determinant
$-7/4$, so aggregate-PSD fails there.
\item[\texttt{two\_keepout\_stage.py}] closes the bracket at
$t^\star=169/320$ on a stage with two active keep-outs whose contacts touch
disjoint coordinates; this is one half of the main theorem's multi-keep-out
generality. Its published certificate is also the paper's floating-point
counterexample object: the same sweep transcribed into IEEE-754 doubles
refuses it, the terminal corner coming out $-4.44\times10^{-16}$ rather than
exactly $0$.
\item[\texttt{overlap\_keepout\_stage.py}] closes the bracket at
$t^\star=5773/3600$ on a stage whose two active contact gradients share
their whole support; this is the other half of the multi-keep-out
generality, and the receipt checks that the verifier's normal-equation
solve recovers $\lambda^\star$ on exactly that stage.
\item[\texttt{hard\_dynamics\_gap.py}] certifies both halves of the paper's
hard-dynamics discussion over $\Q$, on $d{=}1$ instances with one keep-out
and one equality. On the first instance the best constant-multiplier bound
is exactly $1$ against $J^\star=3/2$, a certified gap of $1/2$ whose
diagnosis is an aggregate-PSD failure. On its sibling the aggregate-PSD
analogue holds and the constant-multiplier bracket closes exactly at
$t^\star=J^\star=9/4$.
\end{description}

\paragraph{The pytest pins.}
\begin{description}
\item[\texttt{test\_bracket\_scaling.py}] ships in the bundle and pins the
receipts above: it executes each of the four exact instance receipts and
asserts its published constants, runs \texttt{cw\_anatomy.py}, and checks
\texttt{cert\_metrics.py}'s results-table rows at the headline sizes.
\item[\texttt{test\_riccati}] ships in the Podium library and pins the band
checkers themselves: the band-$H$ sweep on the assembled $H$ and the
bordered sweep on the assembled $M$ are checked against the dense reference
predicate \texttt{is\_psd} of \texttt{podium.verify.barrier}, on randomized
fixed-seed instances plus dedicated singular-$H$ and zero-pivot cases.
\end{description}

\paragraph{The height receipts.} These receipts pin the paper's size
measures and the separations among its named heights.
\begin{description}
\item[\texttt{height\_measure\_check.py}] pins the paper's height and
bit-size definition: the enumerated least height is attained at the
common-denominator form over the lcm of the reduced denominators, per-entry
bit-size obeys $\mathrm{bit}\le2\,\mathrm{ht}$, and a positive integer $D$
read as a rational obeys $\mathrm{bit}(D)=\mathrm{bitlen}(D)+1$.
\item[\texttt{heights\_enumerated.py}] asserts that on the CW-coupled
family the four height-bearing tuples carry four pairwise distinct heights
at every one of the $23$ horizons, in the strict order
$\mathrm{ht}(x^\star)<L<L_M<L_{\mathrm{asm}}$.
\item[\texttt{height\_inversion.py}] exhibits a family admissible under the
paper's constructed-family definition on which the assembled height
$L_{\mathrm{asm}}$ is the \emph{smaller} of the two design-side heights ---
the receipt behind the statement that the definition forces no ordering
between $L$ and $L_{\mathrm{asm}}$.
\item[\texttt{height\_governance.py}] exhibits an admissible family,
closing on all three trusted checkers, with $L_M\le8$ bits against a design
height $L=257$--$258$, so that there the per-stage design height bounds
every entry of the bordered matrix --- the ordering opposite to the CW
family's.
\item[\texttt{claim311\_recheck.py}] extends the certified-value
measurements past the demonstration window by an exhaustive exact sweep of
every horizon $N=2,\dots,6143$: $\den(t^\star)$ stays inside $450$--$526$
bits at every horizon, while $\mathrm{bit}(t^\star)$ keeps the residual
$\log N$ drift.
\item[\texttt{claim258\_const\_lm\_witness.py}] constructs the
constant-$L_M$ witness: an admissible closing family with $L_M=2$ at every
horizon whose largest interior pivot grows without bound, so no function of
$L_M$ alone bounds the sweep's cost.
\item[\texttt{attackb\_scope.py}] constructs the fixed-order witness ---
at order $n{+}1=9$, multipliers of growing denominator raise $L_M$ and the
largest interior pivot together, so no function of the order alone bounds
the cost either. It also separates the multiplier-recovery solve's operand
heights from its formation heights on the demonstration and inverted
families.
\item[\texttt{attack2\_scope.py} and \texttt{attackb\_scope.py}] measure the
recovery solve's operands at the design height: $5$--$6$ bits against $L=6$
on the CW family (\texttt{attackb\_scope.py}), and $259$ bits against
$L=258$ on \texttt{height\_inversion.py}'s inverted family
(\texttt{attack2\_scope.py}).
\item[\texttt{claim235\_thirdcheck.py} and \texttt{claim306\_attack.py}]
jointly pin the two witnesses for why the certified value's bit bound
carries a horizon term and is not stated against $L$ alone: an admissible
closing family with $L_{\mathrm{asm}}=2$, $\mathrm{ht}(x^\star)=1$ and
$t^\star=N$ exactly, and integer goal weights of growing magnitude taking
$\den(t^\star)$ unbounded at fixed $L=1$.
\end{description}

\paragraph{The class-structure and sweep-scope receipts.} These receipts
pin structural facts about the class and the scope of the band sweep.
\begin{description}
\item[\texttt{bandwidth\_scope.py}] pins the scalar bandwidth: the
neighbour-only coupling confines every nonzero entry to half-width
$b=2d-1$, attained whenever some coupling block has a nonzero $(1,d)$
entry, while structured couplings such as $A_k=R_k=I$ sit at half-width
exactly $d$.
\item[\texttt{tridiag\_structure\_scope.py}] pins that block tridiagonality
of $H(\lambda)$ is forced by the class rather than assumed, with a negative
control: the band sweep returns a wrong verdict on a symmetric matrix that
lacks the structure.
\item[\texttt{curvature\_factor\_check.py}] pins the curvature convention
$\nabla^2_{xx}(f_0-\sum_j\lambda_jf_j)=2H(\lambda)$ under the paper's
no-$\tfrac12$ convention, and enumerates which downstream uses of $H$ are
invariant under positive rescaling.
\item[\texttt{band\_lemma\_scope.py}] extends the band-$LDL^\top$ soundness
and zero-pivot checks to general bandwidth, including even $b$, which is
the generality at which the paper states its band lemma.
\item[\texttt{band\_opcount.py}] instruments the shipped band sweep and
pins its $O(nb^2)=O(Nd^3)$ operation count against its loop bounds. It is
the band sweep's count that is measured; the bordered sweep's count is read
off its loop bounds, as the paper's bordered-extension remark states.
\item[\texttt{wdd\_row\_reduction\_check.py}] pins the identity-coupled row
shapes of the paper's diagonal-dominance lemma and the reduction of WDD to
the local test $w_k\ge\lambda_k$; the checks cover the singular boundary,
including the case where the guarded pivot recursion is silent but the
local test is not.
\item[\texttt{riccati\_margin.py}] demonstrates that WDD is sufficient but
not necessary for the PSD verdict: on the identity-coupled chain, PSD
persists below the WDD threshold $c=2$, down to $2\cos(\pi/(N{+}1))$.
\item[\texttt{ya\_scope.py}] pins the boundary of the demonstration
family's growth behaviour: its mild pivot growth and its compact
recomputed value are properties of the family's shared transition matrix.
Swapping the shared CW matrix for the per-stage Yamanaka--Ankersen
transition matrix, with everything else unchanged, reproduces the CW pivot
sequence bit for bit at $e=0$. At the probed eccentricities
$e\in\{0.1,0.3,0.5\}$, with the stage
matrices rationalized to the CW family's realized per-entry fidelity, the
pivots grow at $\approx1000$--$1050$ bits per stage against CW's
$\approx778$, and $\mathrm{bit}(t^\star)$ grows at $\approx1000$ bits per
stage instead of staying horizon-bounded. The penalty is already present
at $e=10^{-4}$: the driver is the transition matrix's per-stage
time-variance, not the eccentricity magnitude.
\end{description}

\paragraph{The boundary receipts.} Five receipts pin the fixed-precision
boundary analysis of the paper's related-work discussion.
\begin{description}
\item[\texttt{claim269\_recheck.py}, \texttt{claim305\_recheck.py} and
\texttt{claim305\_deepchain.py}] pin the positive-margin case: pivots away
from the semidefinite boundary may round, their positive-width enclosures
absorbed by their positive margins.
\item[\texttt{claim305\_recheck2.py}] pins the exactly-computed-zero-pivot
scope of the Cholesky clause: LAPACK's \texttt{dpotrf} declines the
factorization of the singular $H$ with rows $(1,-1)$ and $(-1,1)$. That an
$LDL^\top$ and the symmetric eigensolver are exact there is pinned by
\texttt{attack3\_scope.py} and \texttt{claim269\_recheck.py}.
\item[\texttt{attack3\_scope.py}] pins the floating-point transcription the
paper's conclusion reports --- the IEEE-754 double sweep refusing
\texttt{two\_keepout\_stage.py}'s certificate at a terminal corner of
$-4.44\times10^{-16}$ --- and the zero-pivot branch on which the bordered
and band-$H$ pivot sequences diverge.
\end{description}

\section{Mechanization of the published note's theorems}\label{sec:published-note}
The exact S-procedure bracket and the rational $LDL^\top$ PSD test for
single primitives are established in the published Podium technical note,
\emph{Exact-Rational Certificates in Podium: Constructions, Proofs, and
Prior Art}~\cite{oltean2026certnote}. Its soundness core is mechanized in
\texttt{PodiumNote.lean}, which sits in a \texttt{note/} subdirectory beside
the eighteen top-level files and is unchanged since commit \texttt{91fb1fa}
of the same tree. The regression script of \S\ref{sec:corpus} type-checks
the file as its nineteenth entry, under the same toolchain pin, and surfaces
the axiom audit for \texttt{note\_thm\_sound}, \texttt{note\_thm\_multi} and
\texttt{note\_lem\_schur}. The declarations are stated over any linear
ordered field, like the corpus's, and mechanize the note's Theorem~1, the
lower-bound leg of its Theorem~4 with the accompanying bracket, and the
scalar case of its Schur-complement lemma.

\texttt{note\_thm\_sound} (the note's Theorem~1, the lower bound, at one
constraint): if $\lambda\ge0$ and the minorant condition
$t\le f_0(x)-\lambda f_1(x)$ holds at every point $x$, then $t\le f_0(x)$
at every $x$ with $f_1(x)\ge0$, hence $t\le J^\star$. The minorant
condition is where $M(\lambda,t)\succeq0$ enters: the published note
derives it from PSD of the bordered matrix by the lifting expansion, and
the declaration takes the minorant itself as its hypothesis. Within the
corpus, the same step is carried at general order by
\texttt{general\_sprocedure\_sound}'s form-identity hypothesis
(\S\ref{sec:soundness}).

\texttt{note\_thm\_multi} (the note's Theorem~4, the certified duality gap,
at any number $m$ of constraints): if every $\lambda_i\ge0$ and
$t\le f_0(x)-\sum_i\lambda_if_i(x)$ at every point, then $t\le f_0(x)$ at
every $x$ with all $f_i(x)\ge0$ --- the theorem's opening statement, that
Theorem~1's minorant argument uses nothing of $m=1$. The declaration is
stated at one fixed $(\lambda,t)$; the theorem's supremum characterization
of the Shor value $d^\star$, and the strictness of the gap for $m\ge2$, are
outside it.

\texttt{note\_bracket} (the same theorem's closing observation): a lower
bound $t\le f_0$ over a feasible set, supplied as a hypothesis, together
with a feasible witness $\bar x$, yields $t\le f_0(\bar x)$ and, at every
feasible $x$, $f_0(\bar x)-(f_0(\bar x)-t)\le f_0(x)$ --- the form in which
the width $f_0(\bar x)-t$ of the bracket $[t,f_0(\bar x)]$ bounds the
suboptimality of $\bar x$. Feasibility enters as an abstract predicate, so
the statement is independent of the constraint form.

\texttt{note\_lem\_schur} (the note's Schur-complement lemma, scalar case):
for $a>0$, the bordered matrix
$M=\bigl[\begin{smallmatrix}a&b/2\\b/2&c-t\end{smallmatrix}\bigr]$ is PSD in
the quantifier sense $\forall v,\ 0\le v^\top Mv$ if and only if
$t\le c-b^2/(4a)$; the right-hand side is the dual value
$\inf_x(ax^2+bx+c)$. The equivalence is mechanized in both directions.
The published lemma is stated for a matrix block $A(\lambda)\succ0$ with
$g(\lambda)=c-\tfrac14b^\top A(\lambda)^{-1}b$; the declaration is its
$1\times1$ instance.

\section{Scope: what the theorems take as hypotheses}\label{sec:scope}
The shipped sweeps are not the mechanized object: the corpus proves the
acceptance predicate's equivalence with positive semidefiniteness. The claim
that the concrete band elimination in \texttt{podium.verify.riccati} computes
such a factorization is a refinement hypothesis, taken rather than proved.
The same holds for the dense reference $M$-test of
\texttt{podium.verify.barrier}, which pins the band elimination and against
which every emitted certificate is checked.

Three things are outside the mechanized perimeter: the border-carrying
sweep; the band-restriction equality (band pivots equal dense pivots); and
the operation count. The paper pins the first by test against the dense
$M$-test and reads the count off the shipped loop bounds; the corpus
contains none of the three.

The form identity is discharged per instance:
\texttt{general\_sprocedure\_sound}'s third hypothesis and the capstone's
matching hypothesis are proved by \texttt{ring} for each concrete assembly,
never once in general.

Lagrangian minimality is supplied from outside the corpus:
\texttt{certified\_global\_optimum} takes it as a hypothesis, and the
quadratic-convexity step that produces it in the paper's proof is not
mechanized.

The bridge is compiled separately: \texttt{PosSemidefBridge.lean},
\texttt{LDLTSound.lean} and \texttt{SchurStep.lean} each declare their own
top-level \texttt{IsPSD} under the same unnamespaced name, and the three files
import none of one another. So no compiled artifact links the bridge to the
predicate the acceptance lemmas decide. The three formulas are definitionally
identical over $\Q$; the match is read off by inspection.

The second leg of the paper's diagonal-dominance lemma is the
Schur-complement preservation of weak diagonal dominance, resting on
Carlson--Markham~\cite{carlson1979} together with the limiting argument in
that proof. That leg is outside the corpus; only the Gershgorin leg is
machine-checked.

\bibliographystyle{plain}
\bibliography{references}
\end{document}
